\documentclass[a4paper,11pt]{article}
\usepackage[T1]{fontenc}
\usepackage[utf8]{inputenc}
\usepackage{amsmath}
\usepackage{amssymb}
\usepackage{siunitx}
\usepackage[margin=2.5cm]{geometry}

\title{ATCG I Sheet R05 Solution}
\author{
\large Kamran Moradnejad (3099229) \\
\large Pranav Yadav (50327698) \\
\large Saransh Jain (50262919)
}
\date{May 15, 2026}

\begin{document}
\maketitle

\section*{Assignment 2: Properties of Microfacet BRDFs}

The two renderings in Figure~2 of the sheet show the same rough plane lit by an
environment map. The image that displays a pronounced, elongated brightening of
the reflected environment along the horizon and at grazing viewing angles
corresponds to the \textbf{micro-facet BRDF}. The image whose specular
reflection looks comparatively dull, with a roughly symmetric (cosine-shaped)
falloff and no horizon brightening, corresponds to the \textbf{Phong BRDF}.

\subsubsection*{Why the appearance differs}

The Phong specular term has the form
\[
f^{\text{Phong}}_r(\omega_i,\omega_o)
   = k_s\,\frac{n+2}{2\pi}\,(\hat{r}\cdot\omega_o)^n,
\]
where $\hat r$ is the perfect mirror reflection of $\omega_i$. The lobe only
depends on the angle between $\hat r$ and the viewing direction. The amplitude
of the highlight does not depend on the viewing zenith $\theta_v$ at all, so
when the camera looks across the surface near grazing, Phong produces the same
modest, roughly circular highlight as it does for a downward view.

A Cook--Torrance-style micro-facet BRDF has the form
\[
f^{\text{micro}}_r(\omega_i,\omega_o)
   = \frac{D(h)\,G(\omega_i,\omega_o)\,F(\omega_i\!\cdot\! h)}{4\,(n\cdot\omega_i)(n\cdot\omega_o)}.
\]
Two ingredients of this formula are responsible for the visible difference:
\begin{itemize}
    \item The denominator carries a factor $1/(n\cdot\omega_o)$, which goes to
    infinity as the view approaches grazing. Smith shadowing $G$ also vanishes
    at grazing, but at a slower rate than the cosine in the denominator, so the
    BRDF value diverges as $\theta_v\!\to\!\pi/2$. This is the well known
    ``horizon brightening'' of micro-facet models.
    \item Schlick's Fresnel term grows from $F_0$ at normal incidence to
    essentially~$1$ at grazing. Together with the geometry term it makes the
    reflectance at grazing visibly bright and stretched along the view
    direction.
\end{itemize}
Phong contains neither of these mechanisms, which is why its horizon stays dim
and its highlight stays symmetric. The asymmetric, stretched grazing
reflections in the micro-facet image are therefore an immediate consequence of
the $DGF/(4(n\!\cdot\!\omega_i)(n\!\cdot\!\omega_o))$ structure of the
micro-facet BRDF combined with the angle dependence of $F$ and $G$.

\section*{Assignment 3: Properties of a Heitz Surface}

The Heitz surface is the random height field
\[
h(x,y) = \sqrt{\frac{2}{N}}\sum_{i=1}^{N}\cos(\Phi_i + x f_{ix} + y f_{iy}),
\]
where the phases $\Phi_i \sim \mathcal{U}[0,2\pi]$ are i.i.d.\ and independent
of the frequencies $(f_{ix},f_{iy})$, which are themselves i.i.d.\ samples of a
spectral density. We fix a query point $(x,y)$ throughout and write
\(\Psi_i := \Phi_i + x f_{ix} + y f_{iy}\). Conditional on
$(f_{ix},f_{iy})$, $\Psi_i$ is still uniform on $[0,2\pi]$ modulo $2\pi$, hence
$\sin\Psi_i$ and $\cos\Psi_i$ are zero-mean with second moment $1/2$ and
mean-zero correlation $E[\sin\Psi_i\cos\Psi_i]=0$.

\subsection*{a) Limit of the height distribution}

Let
\[
X_i := \sqrt{\frac{2}{N}}\,\cos\Psi_i,
\qquad i=1,\dots,N,
\]
so that $h(x,y)=\sum_{i=1}^N X_i$. The $X_i$ are i.i.d.\ (their joint
distribution is invariant under permutations because both $\Phi_i$ and
$(f_{ix},f_{iy})$ are i.i.d.\ across $i$). Using the law of total expectation,
\[
E[X_i] = \sqrt{\frac{2}{N}}\,E\!\left[E[\cos\Psi_i\,|\,f_{ix},f_{iy}]\right]
       = \sqrt{\frac{2}{N}}\cdot 0 = 0,
\]
\[
E[X_i^2] = \frac{2}{N}\,E[\cos^2\Psi_i]
        = \frac{2}{N}\cdot\frac{1}{2} = \frac{1}{N}.
\]
Therefore
\[
E[h] = \sum_{i=1}^N E[X_i] = 0,
\qquad
\operatorname{Var}(h) = \sum_{i=1}^N E[X_i^2] = N\cdot\frac{1}{N} = 1.
\]
The summands are i.i.d.\ with mean $0$, variance $1/N$ and are uniformly
bounded by $\sqrt{2/N}$, so the Lindeberg condition is trivially satisfied. By
the classical central limit theorem
\[
h(x,y) \;\xrightarrow[N\to\infty]{d}\; \mathcal{N}(0,1).
\]

\subsection*{b) Limit of the slope distribution}

Differentiating term by term,
\[
\frac{\partial h}{\partial x}
   = -\sqrt{\frac{2}{N}}\sum_{i=1}^N f_{ix}\sin\Psi_i,
\qquad
\frac{\partial h}{\partial y}
   = -\sqrt{\frac{2}{N}}\sum_{i=1}^N f_{iy}\sin\Psi_i.
\]
Set
\[
Y_i := -\sqrt{\frac{2}{N}}\,f_{ix}\sin\Psi_i.
\]
$f_{ix}$ is independent of $\Psi_i$ (conditioning on $f_{ix}$ leaves the
distribution of $\Psi_i$ uniform). Hence
\[
E[Y_i] = -\sqrt{\frac{2}{N}}\,E[f_{ix}]\,E[\sin\Psi_i] = 0,
\]
\[
E[Y_i^2] = \frac{2}{N}\,E[f_{ix}^2]\,E[\sin^2\Psi_i]
        = \frac{2}{N}\,E[f_{ix}^2]\cdot\frac{1}{2}
        = \frac{E[f_{ix}^2]}{N}.
\]
For the Heitz construction the spectrum of frequencies is chosen so that the
second moment of $f_{ix}$ equals
\[
E[f_{ix}^2] = \frac{\alpha_x^2}{2},
\]
and analogously $E[f_{iy}^2] = \alpha_y^2/2$. (This is exactly the property
that makes the limiting slope distribution carry the standard roughness
parameters $\alpha_x,\alpha_y$.) Therefore
\[
\operatorname{Var}\!\left(\frac{\partial h}{\partial x}\right)
   = N\cdot\frac{E[f_{ix}^2]}{N}
   = \frac{\alpha_x^2}{2},
\qquad
\operatorname{Var}\!\left(\frac{\partial h}{\partial y}\right)
   = \frac{\alpha_y^2}{2}.
\]
The summands $Y_i$ are i.i.d.; assuming $f_{ix}$ has a finite second moment
(true by the construction above), Lindeberg's condition holds and the CLT
yields
\[
\frac{\partial h}{\partial x}
   \;\xrightarrow[N\to\infty]{d}\;
   \mathcal{N}\!\left(0,\,\frac{\alpha_x^2}{2}\right),
\qquad
\frac{\partial h}{\partial y}
   \;\xrightarrow[N\to\infty]{d}\;
   \mathcal{N}\!\left(0,\,\frac{\alpha_y^2}{2}\right).
\]
A direct application of the Cramér--Wold device extends this to joint
convergence of the gradient $(\partial_x h,\partial_y h)$ to a two-dimensional
Gaussian with mean $0$ and covariance
\(
\operatorname{diag}(\alpha_x^2/2,\;\alpha_y^2/2)
\)
(the cross terms vanish because
$E[\sin\Psi_i\sin\Psi_j]=0$ for $i\neq j$ and
$E[f_{ix}f_{iy}]=0$ when the spectrum is separable, as is the case here).

\end{document}
